% SHARD-ID: AQM-14-ARITHMETIC-QUANTUM-FIELDS
% SHARD-TITLE: Arithmetic Quantum Fields
% SHARD-SUMMARY: Defines finite arithmetic quantum fields as Weyl displacements of finite symplectic function spaces.
% SHARD-SUMMARY: Proves the pointwise vector-space and symplectic-form lemmas, including the radical caveat.
% SHARD-SUMMARY: Computes exact \(F_3\) examples for finite sets, vector spaces, and arithmetic varieties.
% SHARD-KEYWORDS: arithmetic quantum field, finite field, Weyl-Heisenberg, function space, symplectic form, finite variety

\section{Arithmetic Quantum Fields}
\label{sec:arithmetic-quantum-fields}

\subsection{Status and Source Anchors}

This section introduces a proposal.  The finite-field Weyl facts are
source-local: Ashikhmin--Knill define \(p^m\)-ary shift/phase operators and
their commutator in \path{n9.tex}, lines 249--313, and
Bombin--Martin-Delgado record the qudit Pauli symplectic commutator in
\path{HomologicalErrorCorrection6.tex}, lines 1206--1259.  The new content
below is the checked construction of field label spaces from finite function
spaces.

\subsection{Lamport Dependency Skeleton}

\[
\begin{array}{rcl}
D1&:&X\text{ is a finite physical space.}\\
D2&:&(V,\omega_V)\text{ is a finite symplectic }F\text{-vector space.}\\
D3&:&E\leq \operatorname{Map}(X,V)\text{ is a chosen }F\text{-linear field space.}\\
D4&:&\Omega_E(\phi,\psi)=\sum_{x\in X}w_x\omega_V(\phi(x),\psi(x)).\\
L1&:&\operatorname{Map}(X,V)\text{ is an }F\text{-vector space pointwise.}\\
L2&:&\Omega_E\text{ is alternating and bilinear.}\\
L3&:&\Omega\text{ is nondegenerate on all maps if }\omega_V\text{ is and }w_x\ne0.\\
L4&:&E\text{ is symplectic iff }\operatorname{rad}(E)=0.\\
L5&:&E_{\mathrm{red}}=E/\operatorname{rad}(E)\text{ is symplectic.}\\
L6&:&W:E_{\mathrm{red}}\to U(\mathcal H)\text{ is a projective Weyl representation.}\\
T1&:&\mathcal A_{X,E}=\operatorname{span}_{\mathbb C}\{W(e)\}\text{ is the field-observable algebra.}
\end{array}
\]
The numbering is used only to check that no later construction is invoked
before its inputs are defined.

\subsection{Definitions D1--D4}

Let \(F=\mathbb F_q\), \(q=p^m\).  A finite physical space is first only a
finite set \(X\).  If \(X\) has additional structure, for example an
\(F\)-vector-space or variety structure, then \(\operatorname{Hom}(X,V)\)
means a specified class of structure-preserving maps.  Without that extra
choice the neutral object is
\[
  \operatorname{Map}(X,V)=\{\phi:X\to V\}.
\]
The proposed arithmetic field label space is not automatically all maps.  It
is a chosen \(F\)-linear subspace
\[
  E\leq \operatorname{Map}(X,V).
\]
Examples include all functions, linear maps, affine maps, and polynomial
functions up to a degree bound, all interpreted after evaluation on the finite
point set.

Fix nonzero weights \(w_x\in F^\times\).  In this shard and in the validation
run all weights are \(w_x=1\).  Define
\[
  \Omega_E(\phi,\psi)
  =
  \sum_{x\in X}w_x\,\omega_V(\phi(x),\psi(x)).
\]

\subsection{Pointwise Linear Algebra L1--L5}

\paragraph{Lemma \(L1\).}
\(\operatorname{Map}(X,V)\) is an \(F\)-vector space under pointwise operations:
\[
  (a\phi+b\psi)(x)=a\phi(x)+b\psi(x).
\]
\emph{Proof.}
The right side lies in \(V\), so it defines a map \(X\to V\).  Associativity,
commutativity of addition, distributivity, scalar identity, and additive
inverses hold after evaluating at each \(x\), where they are exactly the vector
space axioms in \(V\).  Thus the function set is a vector space.  Any subset
closed under these operations is an \(F\)-linear field subspace \(E\).

\paragraph{Lemma \(L2\).}
\(\Omega_E\) is \(F\)-bilinear and alternating.
\emph{Proof.}
For \(a,b\in F\),
\[
\begin{aligned}
\Omega_E(a\phi+b\psi,\eta)
&=\sum_xw_x\omega_V(a\phi(x)+b\psi(x),\eta(x))\\
&=a\Omega_E(\phi,\eta)+b\Omega_E(\psi,\eta).
\end{aligned}
\]
Linearity in the second argument is identical.  Also
\[
  \Omega_E(\phi,\phi)=\sum_xw_x\omega_V(\phi(x),\phi(x))=0
\]
because \(\omega_V\) is alternating.  Hence \(\Omega_E\) is an alternating
bilinear form.

\paragraph{Lemma \(L3\).}
If \(E=\operatorname{Map}(X,V)\), \(\omega_V\) is nondegenerate, and every
\(w_x\ne0\), then \(\Omega_E\) is nondegenerate.
\emph{Proof.}
Let \(\phi\ne0\).  Choose \(x_0\) with \(\phi(x_0)\ne0\).  Since
\(\omega_V\) is nondegenerate, choose \(v\in V\) with
\(\omega_V(\phi(x_0),v)\ne0\).  Define \(\psi(x_0)=w_{x_0}^{-1}v\) and
\(\psi(x)=0\) for \(x\ne x_0\).  Then
\[
  \Omega_E(\phi,\psi)=\omega_V(\phi(x_0),v)\ne0.
\]
Thus no nonzero \(\phi\) is orthogonal to all maps.

\paragraph{Lemma \(L4\).}
For a subspace \(E\), define
\[
  \operatorname{rad}(E)=\{\phi\in E:\Omega_E(\phi,\psi)=0
  \text{ for all }\psi\in E\}.
\]
Then \((E,\Omega_E)\) is symplectic iff \(\operatorname{rad}(E)=0\).
\emph{Proof.}
By definition, nondegeneracy means the only vector pairing to zero with every
element of \(E\) is \(0\).  That set is exactly \(\operatorname{rad}(E)\).
Together with Lemma \(L2\), this is precisely the definition of a symplectic
vector space.

\paragraph{Lemma \(L5\).}
The quotient \(E_{\mathrm{red}}=E/\operatorname{rad}(E)\) is symplectic.
\emph{Proof.}
If \(r\in\operatorname{rad}(E)\), then
\(\Omega_E(\phi+r,\psi)=\Omega_E(\phi,\psi)\), so the form descends to the
quotient.  If \([\phi]\) pairs to zero with every \([\psi]\), then \(\phi\)
lies in the radical, so \([\phi]=0\).  Hence the descended form is
nondegenerate.

\subsection{Weyl-Heisenberg Field Displacements L6}

Choose a symplectic basis of \(E_{\mathrm{red}}\), so
\[
  E_{\mathrm{red}}\cong F^n\oplus F^n,\qquad e=(q,p).
\]
Let \(\chi(t)=\exp(2\pi i\,\operatorname{Tr}_{F/\mathbb F_p}(t)/p)\).  On
\(\mathcal H=\mathbb C^{F^n}\), set
\[
  T_q|a\rangle=|a+q\rangle,\qquad
  R_p|a\rangle=\chi(p\cdot a)|a\rangle,\qquad
  W(q,p)=T_qR_p .
\]
Then
\[
  W(q,p)W(q',p')
  =
  \chi(p\cdot q')\,W(q+q',p+p')
\]
and therefore
\[
  W(e)W(f)=\chi(\Omega(e,f))W(f)W(e).
\]
The operators \(W(e)\) are unitary because \(T_q\) permutes the basis and
\(R_p\) is diagonal with unit-modulus entries.  The multiplication law shows
that \(e\mapsto W(e)\) is a projective unitary representation of the additive
group of \(E_{\mathrm{red}}\), with cocycle
\(c(e,f)=\chi(p\cdot q')\).

\subsection{Definition: Arithmetic Quantum Field}

An arithmetic quantum field datum is
\[
  (X,V,E,\Omega_E,W_E)
\]
where \(X\) is finite, \(V\) is finite symplectic over \(F_q\), \(E\leq
\operatorname{Map}(X,V)\) is a chosen finite field-label space, and \(W_E\) is
the Weyl projective representation of \(E_{\mathrm{red}}\).  The field
displacement operators are
\[
  \mathsf A_{X,E}=\{W_E(e):e\in E_{\mathrm{red}}\}.
\]
The observable algebra of this finite arithmetic quantum field is
\[
  \mathcal A_{X,E}
  =
  \operatorname{span}_{\mathbb C}\mathsf A_{X,E}
  \subseteq \operatorname{End}(\mathcal H).
\]
When \(\dim_F E_{\mathrm{red}}=2n\), the local Weyl source says these
operators form an operator basis, so \(\mathcal A_{X,E}\cong M_{q^n}(\mathbb C)\).

\subsection{Molecular F3 Examples}

For the run we take \(V=F_3^2\) with
\(\omega((q,p),(q',p'))=pq'-qp'\).  A scalar function space
\(\mathcal U\leq\operatorname{Map}(X,F_3)\) gives
\[
  E=\mathcal U q\oplus\mathcal U p,\qquad
  B(f,g)=\sum_{x\in X}f(x)g(x),
\]
\[
  \Omega((q,p),(q',p'))=B(p,q')-B(q,p').
\]
If \(m=\dim\mathcal U\) and \(r=\operatorname{rank}B\), then
\[
  \dim E=2m,\quad \operatorname{rank}\Omega=2r,\quad
  \dim\operatorname{rad}(E)=2(m-r).
\]

\[
\begin{array}{l|c|c|c|c|c}
\text{example} & |X| & \dim\mathcal U & r & \dim\operatorname{rad} & |E_{\rm red}|\\
\hline
\text{two points, all maps} & 2&2&2&0&81\\
\text{three points, all maps} & 3&3&3&0&729\\
F_3,\ \langle t\rangle & 3&1&1&0&9\\
F_3,\ \langle1,t\rangle & 3&2&1&2&9\\
F_3,\ \langle1,t,t^2\rangle & 3&3&3&0&729\\
F_3^2,\ \langle x,y\rangle & 9&2&0&4&1\\
F_3^2,\ \text{all maps} & 9&9&9&0&387420489\\
y=x^2,\ \langle1\rangle & 3&1&0&2&1\\
y=x^2,\ \langle x,y\rangle & 3&2&2&0&81\\
y=x^2,\ \langle1,x,y\rangle & 3&3&3&0&729\\
G_m(F_3),\ \langle1,t\rangle & 2&2&2&0&81
\end{array}
\]

The negative controls are meaningful.  For affine fields on \(F_3\), the
constant function has zero norm because \(1+1+1=0\) in \(F_3\), so constants
lie in the radical.  For \(\langle x,y\rangle\) on \(F_3^2\), every sum
\(\sum x^2\), \(\sum xy\), and \(\sum y^2\) vanishes, so the entire field
space is radical and the reduced Weyl label group is trivial.

For the parabola \(X=\{(0,0),(1,1),(2,1)\}\), the checked basis values are
\[
  1=(1,1,1),\qquad x=(0,1,2),\qquad y=(0,1,1),
\]
with Gram matrix
\[
  \begin{pmatrix}
  0&0&2\\
  0&2&0\\
  2&0&2
  \end{pmatrix},
\]
which has rank \(3\).  Thus degree-one ambient coordinate functions already
give the full nondegenerate three-dimensional scalar function algebra on these
three points.

\subsection{Exact Validation}

The producer is \path{arithmetic_quantum_fields.jl} under
\path{scripts/bridges/}.  It writes the 2026-05-24 arithmetic-quantum-fields
run bundle.  The summary CSV records point counts, scalar Gram ranks,
symplectic ranks, radicals, reduced Weyl label counts, Hilbert dimensions, and
observable-algebra dimensions.  The basis CSV records every scalar basis vector
and every scalar Gram row used in the examples above.
